% =====================================================================
%  DRAFT SECTION -- preliminary measurements (Phase 1 feasibility probe)
%  Goes in Methods, before the algorithm description: it is what the
%  experimental design was derived FROM.
%
%  Numbers are from phase1_timing.py, run 25.09.2026 on the 15-core
%  machine. UPDATE the genotype length once the controller inputs are
%  final (see the \TODO), and re-run the probe if the architecture or
%  SIM_DURATION changes.
% =====================================================================

\subsection{Preliminary measurements}
\label{sec:preliminary}

Unlike body-scoring tasks, where a candidate is evaluated in microseconds,
every evaluation here requires a physics rollout. Two properties of that
rollout had to be established before the experiment could be designed at
all: how long one evaluation takes, which bounds the population size,
generation count and number of independent runs we can afford; and
whether an evaluation is deterministic, which decides whether the
success rate driving the $1/5$ rule is a reliable signal or partly noise.
Neither can be assumed, so we measured both before implementing the EA.

We compiled the gecko body in \texttt{SimpleFlatWorld} and ran ten
headless rollouts of $15$\,s simulated time with randomly drawn
controller weights, timing them two ways: recompiling the world for every
evaluation, as the provided template does, and compiling once and
resetting the simulation state between evaluations. Determinism was
tested by evaluating one fixed weight vector twice on the same compiled
model, and once more on a freshly compiled one, to rule out state leaking
between evaluations. Table~\ref{tab:preliminary} reports the results.

\begin{table}[h]
    \centering
    \small
    \begin{tabularx}{\columnwidth}{l|X}
        Controller inputs (\texttt{len(data.qpos)}) & $15$ \\
        Controller outputs (\texttt{model.nu}) & $8$ \\
        Genotype length & \TODO{168} weights \\
        Time per evaluation (model reused) & $0.168$\,s \\
        Time per evaluation (recompiled) & $0.210$\,s \\
        Speed relative to real time & $89\times$ \\
        Fitness difference over repeated evaluations & $0.0$ \\
    \end{tabularx}
    \caption{Preliminary measurements, from ten rollouts of $15$\,s
    simulated time. Timings are per evaluation on one core.}
    \label{tab:preliminary}
\end{table}

Three findings shaped the rest of the design.

First, at $0.168$\,s per evaluation a full experiment is inexpensive:
three configurations over twenty independent runs of
\TODO{50}~individuals for \TODO{300}~generations amounts to roughly
$\TODO{900}\,000$ evaluations, or some \TODO{4}~hours of wall-clock time
across \TODO{15} parallel processes. Because compute was not the binding
constraint, we spent it on statistical power rather than on scale: we use
\textbf{twenty independent runs per configuration} instead of the five
required. At five runs per group the smallest attainable two-sided
Mann--Whitney $p$-value is $0.0079$, reached only under complete
separation, so the test bottoms out before it can resolve a moderate
effect; at twenty runs per group that floor falls by several orders of
magnitude.

Second, evaluation is \emph{exactly} deterministic: repeated evaluations
of identical weights returned bit-identical fitness, both on a reused and
on a freshly compiled model. This matters for the research question
rather than for convenience. The $1/5$ rule adapts the step size from the
fraction of offspring that improve on their parent, and under a noisy
fitness some of those improvements would be measurement error, biasing
the signal that drives the controller. Determinism removes that concern
entirely, makes repeated rollouts per individual unnecessary, and means
each run is exactly reproducible from its seed.

Third, reusing the compiled model saves only $1.25\times$, which tells us
that almost all of the time is spent stepping the physics rather than
building the model. We adopt the reuse because it is free, but we did not
pursue further optimisation: the remaining cost is irreducible without
shortening the simulation, which would change the task rather than speed
it up.
